按照第 3.6.1 节 DMA 矩阵乘法的块与线程划分，追踪输出位置 \co@(37, 50)@。该例使用 8 × 16 的块网格，每块负责 16 × 16 的输出分块。这与前面标量示例中的 16 × 64 逻辑网格不同，但两者覆盖相同的输出形状。

各块把 128 行划分为 8 个大小为 16 的分块：\co@px = 37 / 16 = 2@，偏移 \co@qx = 37 % 16 = 5@。列方向把 256 划分为 16 个大小为 16 的分块：\co@py = 50 / 16 = 3@，偏移 \co@qy = 50 % 16 = 2@。于是块为 \co@(2, 3)@，线程为 \co@(5, 2)@。

对该线程，K 循环执行 16 次迭代。在 \co@tile_k = 0@ 这次迭代中，\co@dma.copy@ 把 \co@lhs@ 的第 32..47 行、K 的第 0..15 列装入共享的 \co@lhs_load@，把 \co@rhs@ 的 K 第 0..15 行、第 48..63 列装入共享的 \co@rhs_load@。内层 \co@foreach k in [16]@ 把 \co@lhs_load.data.at(5, k) * rhs_load.data.at(k, 2)@（k = 0..15）累加到 \co@output.at(37, 50)@。随后 \co@tile_k = 1@ 装入下一条 K 条带，依此类推。16 次迭代之后，\co@output.at(37, 50)@ 中就是完整的点积——与标量参考版本的答案相同。
